\chapter{Receptury}
\markboth{Receptury}{Receptury}

Wszystkie receptury w tej książce są podane w gramach i w procentach piekarskich. Procent piekarski to masa składnika podzielona przez masę mąki i pomnożona przez sto. Mąka zawsze ma sto procent. Dzięki temu recepturę można przeskalować bez przeliczania od nowa: zmieniasz masę mąki, a reszta składników zachowuje swoje procenty.

Dwie wartości w tabelach warto rozumieć dokładnie. \emph{Nawodnienie} to masa wody podzielona przez masę mąki, wyrażona w procentach. \emph{Masa całkowita} to suma wszystkich składników, bez strat w procesie. Obie są liczone z gramów przez narzędzie, a nie wpisywane. W tej książce nawodnienie mieści się między \val{bread.hydration.min} a \val{bread.hydration.max} procent, a receptur jest \val{breads.total}.

Zakwas traktujemy jako zaczyn o nawodnieniu sto procent: sto gramów zakwasu to pięćdziesiąt gramów mąki i pięćdziesiąt gramów wody. Dlatego w chlebie pszennym na zakwasie 100 g zakwasu dokłada do mąki 50 g, a do wody 50 g. Stąd 450 + 50 = 500 g mąki, 300 + 50 = 350 g wody i nawodnienie 350 / 500 = 70 procent. Ten rachunek powtarzaj przy każdym chlebie, zanim sięgniesz po wagę.

Sól w większości receptur to około dwóch procent mąki: 10 g na 500 g. Drożdży w tabeli jest od 1 do 5 g na 500 g mąki, czyli od 0,2 do 1 procenta. Zakwas w chlebach żytnio-pszennym to 150 g, czyli 30 procent mąki. Przelicz składniki na procenty, zanim zmienisz którąkolwiek liczbę.

\subsection*{Składniki i masy}

\generated{recipes}

\subsection*{Czasy i temperatury}

\generated{recipe-timing}

Kolumna \emph{zimno} to godziny fermentacji w lodówce. Kolumna \emph{pokojowa} to godziny w temperaturze pokojowej. Kolumny \emph{piec} i \emph{pieczenie} to temperatura pieca i czas w nim. Wszystkie te godziny w planie są czasem biernym, dlatego w Części III zobaczysz je jako harmonogramy.

\section*{Metoda wspólna dla chlebów}

Metoda jest jedna dla wszystkich receptur. Różnią się składnikami i czasami z tabel, a nie kolejnością pracy.

\begin{enumerate}
\item \textbf{Zważ wszystko w gramach.} Waga kuchenna z dokładnością do jednego grama wystarczy. Mąki nie odmierzaj szklankami.
\item \textbf{Przygotuj zakwas albo drożdże.} Zakwas o nawodnieniu sto procent odśwież przed użyciem i wykorzystaj go wtedy, gdy jest wyraźnie wyrośnięty. Ocena dotykiem i wyglądem jest ważniejsza niż zegarek.
\item \textbf{Wymieszaj mąkę z wodą i odpocznij.} Krótki odpoczynek mąki z wodą trwa \verify{zwykle kilkadziesiąt minut}, zanim doda się resztę.
\item \textbf{Dodaj sól, zakwas lub drożdże i wyrabiaj.} Wyrabiaj do uzyskania ciasta gładkiego i sprężystego. Ten stan oceniasz dotykiem, nie zegarkiem.
\item \textbf{Fermentuj w temperaturze pokojowej.} Przez czas z kolumny \emph{pokojowa}. Ciasto powinno wyraźnie zwiększyć objętość. Kolumna mówi, ile godzin, nie ile razy ma urosnąć.
\item \textbf{Uformuj bochenek i daj mu odpocząć.} Formuj tak, żeby na powierzchni powstało napięcie, a potem przykryj.
\item \textbf{Fermentuj na zimno.} W lodówce przez czas z kolumny \emph{zimno}, szczelnie przykryte. Lodówka powinna utrzymywać temperaturę \verify{poniżej 5 °C}, bo tylko wtedy fermentacja zwalnia do poziomu, który planujesz.
\item \textbf{Nagrzej piekarnik wcześniej.} Do temperatury z tabeli. Nagrzewanie zwykle trwa \verify{kilkanaście minut}.
\item \textbf{Piecz w temperaturze i czasie z tabeli.} Po upływie czasu sprawdź wewnętrzną temperaturę termometrem, jeśli go masz. Wnętrze chleba pszennego zwykle osiąga \verify{około 96 °C}.
\item \textbf{Ostudź na kratce.} Krojenie ciepłego chleba psuje strukturę miękiszu. Studzenie zajmuje zwykle \verify{co najmniej godzinę}.
\end{enumerate}

\begin{piekarnik}
Piekarniki różnią się od nastawionej temperatury, \verify{czasem o kilkanaście stopni}. Kup termometr do piekarnika, zmierz, ile pokazuje wnętrze przy nastawie z tabeli, i ustaw czas według tego pomiaru, nie według pokrętła. Pierwszy bochenek traktuj jako pomiar, nie jako dowód. Zapisz w swoim szablonie z Dodatku C, jak wypadł, i co zmienisz przy następnym.
\end{piekarnik}

Receptura nie jest prawem. Jeśli mąka, zakwas albo kuchnia zachowują się inaczej niż w tabeli, zmień czas, a nie tabelę. Tabela pokazuje punkt wyjścia i sposób myślenia o cieście, a nie wynik, który musisz osiągnąć.
